% Compile with latexmk (uses ~/.latexmkrc: LuaLaTeX, output in build/)
% Read-aloud teaching script for Week 9 (companion to beamer.tex, follows it
% frame by frame). Sessions 1 and 2 are held on the SAME DAY, back to back
% (2 x 90 min with a short break in between).
% Exercises are numbered by week (9.1, 9.2, ...); finance examples share the
% numbering. Each one lives in exercises/NN-slug.tex.

\documentclass[11pt,a4paper,parskip=half]{scrartcl}
\usepackage{amsmath}
\usepackage[stixtwo]{fontsetup}
\usepackage[english]{babel}
\usepackage{booktabs}
\usepackage{enumitem}
\usepackage{tikz}
\usepackage[most]{tcolorbox}
\usepackage[a4paper,margin=1in,marginparwidth=1.9cm,marginparsep=2mm]{geometry}
\usepackage[hidelinks]{hyperref}

\definecolor{navy}{RGB}{24,48,75}
\definecolor{teal}{RGB}{0,120,125}
\definecolor{gold}{RGB}{190,130,20}
\definecolor{plum}{RGB}{110,60,120}

\addtokomafont{disposition}{\color{navy}}
\setkomafont{section}{\Large\bfseries\color{navy}}

% Contact
\newcommand{\myemail}{tiborsandor2001@gmail.com}

\newcommand{\Om}{\Omega}
% Expectation, variance and standard deviation, as in Ross: E[X], Var(X), SD(X)
\newcommand{\E}{\operatorname{E}}
\newcommand{\Var}{\operatorname{Var}}
\newcommand{\SD}{\operatorname{SD}}

% Exercise statement (what students see on the slide)
\newtcolorbox{task}[1]{enhanced,breakable,colback=navy!4,colframe=navy,
  fonttitle=\bfseries,title={#1}}
% Finance example
\newtcolorbox{fin}[1]{enhanced,breakable,colback=teal!5,colframe=teal,
  fonttitle=\bfseries,title={#1}}
% Teaching cue, not read aloud
\newtcolorbox{cue}{enhanced,breakable,colback=gold!8,colframe=gold,
  boxrule=.4pt,left=4pt,right=4pt,top=2pt,bottom=2pt,fontupper=\small\itshape}
% Background for the lecturer (finance theory), NOT read aloud
\newtcolorbox{bg}[1]{enhanced,breakable,colback=plum!4,colframe=plum,
  boxrule=.6pt,fonttitle=\bfseries\small,
  title={Background for you (not read aloud): #1},fontupper=\small,
  before upper={\setlength{\parskip}{4pt}}}
% Print version (script-print.tex defines \scriptprint): the background
% boxes are dropped entirely.
\ifdefined\scriptprint
  \RenewDocumentEnvironment{bg}{m +b}{}{}
\fi
\newcommand{\terms}{\par\textbf{\color{plum}Key terms.}\ }
\newcommand{\sq}[1]{\par\textbf{\color{plum}Student may ask:}\ \emph{#1}\par}
\newcommand{\ans}{\textbf{Answer:}\ }

\newcommand{\say}{\textbf{\color{teal}Say:}\ }
\newcommand{\board}{\textbf{\color{gold}Board:}\ }
\newcommand{\ask}{\textbf{\color{navy}Ask the class:}\ }
\newcommand{\slide}[1]{\marginpar{\raggedright\footnotesize\color{gray}slide #1}}

% Small Venn diagram of two events E, F in Omega (same as in the deck)
\newcommand{\venn}[1]{%
  \begin{tikzpicture}[scale=.5,baseline=(O.base)]
    \node (O) at (.5,0) {\phantom{x}};
    #1
    \draw[thick] (-1.5,-1.15) rectangle (2.5,1.15);
    \draw[thick,navy] (0,0) circle (.85);
    \draw[thick,navy] (1,0) circle (.85);
    \node at (-.35,0) {\scriptsize $E$};
    \node at (1.35,0) {\scriptsize $F$};
    \node at (2.15,.8) {\scriptsize $\Omega$};
  \end{tikzpicture}}
\newcommand{\shadeall}[1]{\fill[#1] (-1.5,-1.15) rectangle (2.5,1.15);}

% The exercises (one file each in exercises/) and \scriptexercise
% Exercises of Week 9, shared by beamer.tex, the study guide, the script and
% exercises.tex. Every exercise lives in a file of its own,
% exercises/NN-slug.tex, numbered 9.NN (see exercises.lua for the docstring).
% The body of an exercise file:
%
%   \begin{problem} ... \end{problem}      the statement, used everywhere;
%                                          inside it, \begin{context} ...
%                                          \end{context} is background that
%                                          the slides leave out
%   \begin{slideproblem} ... \end{slideproblem}
%                                          (optional) a shorter statement
%                                          for the slides
%   \begin{slides} ... \end{slides}        one frame of the deck: \exstatement
%                                          places the statement; repeat the
%                                          environment for more frames
%   \begin{solution} ... \end{solution}    the detailed solution (study
%                                          guide, exercise PDFs, website)
%
% Commands for the documents (the argument is the slug: 01-license-plates.tex
% has the slug license-plates):
%
%   \exref{slug}                 the number, e.g. 9.1, for "Exercise~\exref{...}"
%   \extitle{slug}               the title from the docstring
%   \exerciseframe[opts]{slug}   beamer: the frames of the exercise
%   \studyexercise{slug}         study guide: heading, statement, solution
%   \scriptexercise{slug}        script: the statement box, as on the slide
%   \allexercises                exercises.tex: every exercise, in order; on
%                                the website statement and solution with
%                                markers where the pipeline splits the HTML
%                                (\collectionexercise), in the PDF plain:
%                                each statement, followed by its solution if
%                                \ifsolutions (exercises-solutions.tex)

\def\exweek{9}

\directlua{exercises = dofile("exercises/exercises.lua")
  exercises.scan("exercises", "\exweek")}

\ExplSyntaxOn

\msg_new:nnn { exercises } { unknown } { Unknown~exercise~`#1'. }

\cs_new_protected:Npn \__ex_missing:n #1
  { \msg_error:nnn { exercises } { unknown } {#1} ?? }
\cs_new:Npn \exref #1
  {
    \cs_if_exist:cTF { exnum@#1 }
      { \use:c { exnum@#1 } } { \__ex_missing:n {#1} }
  }
\cs_new:Npn \extitle #1
  {
    \cs_if_exist:cTF { extitle@#1 }
      { \use:c { extitle@#1 } } { \__ex_missing:n {#1} }
  }

% ---------------------------------------------------------------------------
% Loading an exercise: its parts are stored, nothing is typeset. The slides
% typeset the stored parts; the other documents read the part they need from
% the file again (\__ex_input_part:nn), as stored tables do not survive lwarp.

\tl_new:N \g__ex_current_tl
\tl_new:N \g__ex_problem_tl
\tl_new:N \g__ex_slideproblem_tl
\tl_new:N \g__ex_solution_tl
\seq_new:N \g__ex_slides_seq
\box_new:N \l__ex_scratch_box

\cs_new_protected:Npn \__ex_load:n #1
  {
    \cs_if_exist:cF { exfile@#1 } { \__ex_missing:n {#1} }
    \tl_gset:Nn \g__ex_current_tl {#1}
    \tl_gclear:N \g__ex_problem_tl
    \tl_gclear:N \g__ex_slideproblem_tl
    \tl_gclear:N \g__ex_solution_tl
    \seq_gclear:N \g__ex_slides_seq
    \vbox_set:Nn \l__ex_scratch_box
      {
        \DeclareDocumentEnvironment { problem } { +b }
          { \tl_gset:Nn \g__ex_problem_tl {##1} } { }
        \DeclareDocumentEnvironment { slideproblem } { +b }
          { \tl_gset:Nn \g__ex_slideproblem_tl {##1} } { }
        \DeclareDocumentEnvironment { slides } { +b }
          { \seq_gput_right:Nn \g__ex_slides_seq {##1} } { }
        \DeclareDocumentEnvironment { solution } { +b }
          { \tl_gset:Nn \g__ex_solution_tl {##1} } { }
        \exp_args:Nv \input { exfile@#1 }
      }
  }

% Typesets one part (problem, slideproblem, solution) of an exercise straight
% from its file; the other parts are skipped.
\cs_new_protected:Npn \__ex_input_part:nn #1#2
  {
    \group_begin:
      \clist_map_inline:nn { problem, slideproblem, slides, solution }
        {
          \str_if_eq:nnTF {##1} {#2}
            { \DeclareDocumentEnvironment {##1} { } { } { } }
            { \DeclareDocumentEnvironment {##1} { +b } { } { } }
        }
      \exp_args:Nv \input { exfile@#1 }
    \group_end:
  }

% Exercise 9.1, or Exercise 9.2 · Finance
\cs_new:Npn \__ex_heading:
  {
    Exercise~\exref { \g__ex_current_tl }
    \cs_if_exist:cT { exfinance@ \g__ex_current_tl } { ~·~Finance }
  }
\prg_new_conditional:Npnn \__ex_if_finance: { T, F, TF }
  {
    \cs_if_exist:cTF { exfinance@ \g__ex_current_tl }
      { \prg_return_true: } { \prg_return_false: }
  }

% ---------------------------------------------------------------------------
% Beamer

% The finance context is printed in the study guide, the exercise PDFs and on
% the website; the
% slides and the script (which shows what the slide shows) leave it out.
\bool_new:N \l__ex_context_bool
\bool_set_true:N \l__ex_context_bool
\DeclareDocumentEnvironment { context } { }
  {
    \bool_if:NTF \l__ex_context_bool
      { \par \textit { Finance~context. } ~ }
      { \__ex_gobble_context:w }
  }
  { \bool_if:NT \l__ex_context_bool { \par } }
\cs_new_protected:Npn \__ex_gobble_context:w #1 \end #2
  { \str_if_eq:nnTF {#2} { context } { \end { context } } { \__ex_gobble_context:w } }

\IfClassLoadedTF { beamer }
  {
    \bool_set_false:N \l__ex_context_bool

    \cs_new:Npn \__ex_block: { \__ex_if_finance:TF { exampleblock } { block } }
    \NewDocumentCommand \exstatement { }
      {
        \begin { \__ex_block: } { \__ex_heading: }
          \tl_if_empty:NTF \g__ex_slideproblem_tl
            { \g__ex_problem_tl } { \g__ex_slideproblem_tl }
        \end { \__ex_block: }
      }

    \tl_new:N \l__ex_frames_tl
    \NewDocumentCommand \exerciseframe { O{t} m }
      {
        \__ex_load:n {#2}
        \tl_clear:N \l__ex_frames_tl
        \seq_map_indexed_inline:Nn \g__ex_slides_seq
          {
            \tl_put_right:Nn \l__ex_frames_tl { \begin{frame}[#1] }
            \tl_put_right:Ne \l__ex_frames_tl
              {
                {
                  \exp_not:N \extitle {#2}
                  \int_compare:nNnT {##1} > { 1 } { \c_space_tl (cont.) }
                }
              }
            \tl_put_right:Nn \l__ex_frames_tl { ##2 \end{frame} }
          }
        \tl_use:N \l__ex_frames_tl
      }
  }
  {
    \cs_new:Npn \__ex_box: { \__ex_if_finance:TF { fin } { task } }
    \cs_new_protected:Npn \__ex_statement_box:n #1
      {
        \begin { \__ex_box: } { \__ex_heading: }
          \__ex_input_part:nn { \g__ex_current_tl } {#1}
        \end { \__ex_box: }
      }

    \NewDocumentCommand \studyexercise { m }
      {
        \__ex_load:n {#1}
        \subsection* { Exercise~\exref{#1}~---~\extitle{#1} }
        \addcontentsline { toc } { subsection }
          { Exercise~\exref{#1}~---~\extitle{#1} }
        \__ex_statement_box:n { problem }
        \soltitle \__ex_input_part:nn {#1} { solution }
      }

    \NewDocumentCommand \scriptexercise { m }
      {
        \__ex_load:n {#1}
        \group_begin:
          \bool_set_false:N \l__ex_context_bool
          \tl_if_empty:NTF \g__ex_slideproblem_tl
            { \__ex_statement_box:n { problem } }
            { \__ex_statement_box:n { slideproblem } }
        \group_end:
      }

    % The website (lwarp): statement and solution of each exercise, between
    % markers where the pipeline splits the HTML.
    \NewDocumentCommand \collectionexercise { m }
      {
        \__ex_load:n {#1}
        \par [[exercise:#1:problem]] \par
        \__ex_input_part:nn {#1} { problem }
        \par [[exercise:#1:solution]] \par
        \__ex_input_part:nn {#1} { solution }
        \par [[exercise:#1:end]] \par
      }
    % The PDF: plain, each statement, followed by its solution in the
    % version with solutions (\ifsolutions). The statement is set off like a
    % quote in Markdown, by a bar on the left (quote=<color>, also used by
    % exercises.tex) in the color of its box in the study guide: navy, or teal
    % for a finance exercise.
    \tcbset
      {
        quote/.style =
          {
            enhanced, breakable, frame~hidden, interior~hidden,
            sharp~corners, boxrule = 0pt, borderline~west = {2pt}{0pt}{#1},
            left = 8pt, right = 0pt, top = 1pt, bottom = 1pt,
            before~skip = 6pt, after~skip = 6pt,
            % the document's paragraph spacing (parskip=half)
            before~upper = { \setlength { \parskip } { .5\baselineskip } }
          }
      }
    \NewDocumentCommand \plainexercise { m }
      {
        \__ex_load:n {#1}
        \subsection* { Exercise~\exref{#1}~---~\extitle{#1} }
        \addcontentsline { toc } { subsection }
          { Exercise~\exref{#1}~---~\extitle{#1} }
        \__ex_if_finance:TF
          { \begin{tcolorbox}[quote=teal] }
          { \begin{tcolorbox}[quote=navy] }
          \__ex_input_part:nn {#1} { problem }
        \end{tcolorbox}
        \legacy_if:nT { solutions }
          {
            \soltitle
            \__ex_input_part:nn {#1} { solution }
          }
      }
    \NewDocumentCommand \allexercises { }
      {
        \IfPackageLoadedTF { lwarp }
          { \directlua { exercises.each("collectionexercise") } }
          { \directlua { exercises.each("plainexercise") } }
      }
  }

\ExplSyntaxOff


% PDF properties and the date on the title page
% The properties of every PDF of Week 9: author (with email address), title,
% subject, keywords and the date the document last changed, in the PDF
% (hyperref) and on its title page (\date), and the header and footer of the
% printed documents (below). Each document loads this file after hyperref and
% then calls \documentinfo{<what it is>}, e.g. \documentinfo{Slides}.
%
% The date: the website build (scripts/build-content.mjs) puts
% \def\lastmodified{YYYY-MM-DD} in front of the document, the date of the last
% commit that changed its sources. Compiled locally, it is today.

\newcommand{\courseinfo}{Mathematics for Finance}
\newcommand{\weeklabel}{Week 9}
\newcommand{\weekinfo}{\weeklabel: Probability, conditional probability, independence, random variables}
\newcommand{\authorname}{Tibor SÁNDOR}
\providecommand{\myemail}{tiborsandor2001@gmail.com}

\makeatletter
\ifdefined\lastmodified\else
	\edef\lastmodified{\the\year-\two@digits{\the\month}-\two@digits{\the\day}}
\fi
\def\di@split#1-#2-#3\@nil{%
	\def\di@year{#1}\def\di@month{#2}\def\di@day{#3}}
\expandafter\di@split\lastmodified\@nil
% November 11, 2026
\edef\lastmodifiedtext{%
	\ifcase\numexpr\di@month\relax\or January\or February\or March\or April\or
		May\or June\or July\or August\or September\or October\or November\or
		December\fi\space\number\di@day, \di@year}
% D:20261111000000, the date format of PDF metadata
\edef\di@pdfdate{D:\di@year\di@month\di@day000000}

\newcommand{\documentinfo}[1]{%
	\hypersetup{
		pdftitle={\courseinfo, \weekinfo\space -- #1},
		pdfauthor={\authorname\space<\myemail>},
		pdfsubject={\courseinfo, \weekinfo},
		pdfkeywords={probability, combinatorics, mathematics for finance,
				Corvinus University of Budapest},
		pdfcreationdate={\di@pdfdate},
		pdfmoddate={\di@pdfdate},
		pdflang={en-US},
	}%
	\date{Last updated: \lastmodifiedtext}%
	\gdef\documentkind{#1}%
}

% Header and footer of the printed documents (KOMA-Script, scrlayer-scrpage;
% the slides have beamer's footline, the website neither): the course, the
% week and the document on the left of the header (the week's topic is on the
% title page), the university on the right; in the footer the author and the
% page as <current>/<total>. The title page gets only the footer.
\newcommand{\documentkind}{}
\IfClassLoadedTF{beamer}{}{\IfPackageLoadedTF{lwarp}{}{%
	\usepackage{lastpage}%
	\usepackage[headsepline=1.5pt,footsepline=1.5pt,plainfootsepline]
		{scrlayer-scrpage}%
	\clearpairofpagestyles
	\setkomafont{pageheadfoot}{\small\sffamily\color{navy}}%
	\setkomafont{pagenumber}{}%
	\setkomafont{headsepline}{\color{navy!35}}%
	\setkomafont{footsepline}{\color{navy!35}}%
	\ihead{\courseinfo\ \textperiodcentered\ \weeklabel\ \textperiodcentered\ 
		\documentkind}%
	\ohead{Corvinus University of Budapest}%
	\ifoot*{\authorname\ \textperiodcentered\ \myemail}%
	\ofoot*{\pagemark/\pageref*{LastPage}}%
	\pagestyle{scrheadings}%
	% Extra stretch for breaking lines with long inline formulas.
	\setlength{\emergencystretch}{1.5em}%
}}
\makeatother

\documentinfo{Teaching script\ifdefined\scriptprint\space(print version)\fi}

\title{Mathematics for Finance}
\subtitle{Week 9 --- Probability, conditional probability, independence,
  random variables\\[2pt]
  \normalsize Teaching script for two 90-minute practical classes held on the
  same day, back to back\ifdefined\scriptprint\\[2pt]
  (print version, without the background boxes)\fi}
\author{Tibor SÁNDOR\\[2pt]
  \normalsize\href{mailto:\myemail}{\texttt{\myemail}}}

\begin{document}
\maketitle

\begin{cue}
	How to use this script. Plain text is meant to be read out loud or
	paraphrased. \textbf{Say} marks key sentences; \textbf{Ask} marks questions
	to put to the class (wait for answers before revealing the solution);
	\textbf{Board} marks what to write. Gold boxes like this one are short notes
	for you only.\ifdefined\scriptprint\else\ \textbf{Purple ``Background for
			you'' boxes} explain the finance behind the examples (what the terms mean,
		why the example matters, what students may ask); they are \emph{not} meant
		to be read out, but they let you answer follow-up questions with
		confidence.\fi\ The script follows the
	deck frame by frame; the slide numbers in the margin are the frame numbers
	in the footline of the deck. Exercise numbers match the slides.
	This week the pace is deliberately slow: every new idea is introduced from
	scratch with a tiny example, and before each exercise the students try it
	themselves (the script says for how long); then you walk through the
	solution in the order in which the slide reveals it.
	Timings are approximate; items marked \emph{(optional)} can be dropped if
	you run late. Only the exercises of the slides are in this script.
\end{cue}

\subsection*{Timing overview}

Both sessions are on the same day: Session 1 (90 min), a short break,
Session 2 (90 min). The timings leave about 6 minutes of reserve in both
sessions.

\begin{center}
	\small
	\begin{tabular}{@{}llrr@{}}
		\toprule
		Slides & Session 1 & min & cum. \\\midrule
		1--3   & Welcome, recap of week 8, plan for today          & 4  & 4  \\
		4      & From counting to probability                      & 7  & 11 \\
		5--6   & Exercise~\exref{two-exams}                        & 9  & 20 \\
		7      & Equally likely outcomes                           & 3  & 23 \\
		8--9   & Exercise~\exref{five-cards}                       & 8  & 31 \\
		10     & Conditional probability                           & 6  & 37 \\
		11--12 & Exercise~\exref{exams-conditional}                & 8  & 45 \\
		13--14 & Exercise~\exref{actuarial-exams}                  & 8  & 53 \\
		15     & Total probability and Bayes' formula              & 5  & 58 \\
		16--17 & Exercise~\exref{early-warning}                    & 11 & 69 \\
		18     & Odds                                              & 4  & 73 \\
		19--20 & Exercise~\exref{early-warning-odds}               & 8  & 81 \\
		21     & Summary                                           & 3  & 84 \\
		       & Reserve (questions, students working alone)       & 6  & 90 \\\midrule
		       & \textbf{Total}                                    &    & \textbf{90} \\\bottomrule
	\end{tabular}
\end{center}

\begin{center}
	\small
	\begin{tabular}{@{}llrr@{}}
		\toprule
		Slides & Session 2 & min & cum. \\\midrule
		22--23 & Session 2 title, recap of Session 1               & 4  & 4  \\
		24     & The generalized Bayes formula                     & 5  & 9  \\
		25--27 & Exercise~\exref{default-by-rating}                & 10 & 19 \\
		28     & Independent events                                & 6  & 25 \\
		29--30 & Exercise~\exref{pairwise-not-mutual}              & 7  & 32 \\
		31--33 & Exercise~\exref{bond-defaults}                    & 9  & 41 \\
		34     & Discrete random variables, distribution function  & 6  & 47 \\
		35--37 & Exercise~\exref{urn-game}                         & 12 & 59 \\
		38     & Expected value, variance, standard deviation      & 6  & 65 \\
		39--40 & Exercise~\exref{call-option}                      & 9  & 74 \\
		41--42 & Exercise~\exref{trading-strategy}                 & 7  & 81 \\
		43     & Summary                                           & 3  & 84 \\
		       & Reserve (questions, students working alone)       & 6  & 90 \\\midrule
		       & \textbf{Total}                                    &    & \textbf{90} \\\bottomrule
	\end{tabular}
\end{center}

\tableofcontents

%======================================================================
\section{Session 1: Probability, conditional probability, Bayes' formula}
%======================================================================

\subsection{Welcome, recap and plan for today (4 min)}\slide{1--3}

\begin{cue}
	Slide 1 (title) is on the screen while they settle down. No introduction
	of yourself is needed: they know you from last week.
\end{cue}

Good morning, everyone, and welcome back. I hope you all survived the
counting of last week.

\say Let me remind you where we are. Last week we learned to count: the
basic principle of counting, permutations, combinations, the multinomial
coefficient and stars and bars. At the end we met the language of
probability: the sample space $\Om$, the set of all possible outcomes, and
events, which are subsets of $\Om$. We combined events with ``or'', ``and''
and ``not'', we used De Morgan's laws, and we filled Venn diagrams from the
inside out.

Today we put numbers on events: we say \emph{how likely} they are. That is
what probability is about.

\begin{cue}
	Slide 2 (Plan for today).
\end{cue}

\say Today is again a long day: two sessions of 90 minutes, back to back,
with a short break in between.

In Session 1 we learn the basic rules of probability: how to compute the
probability of ``or'', ``and'' and ``not''. Then we see what happens when all
outcomes are equally likely: there, last week's counting gives the
probabilities. After that comes the most important idea of the day:
\emph{conditional probability}, the probability of an event when we already
know something. It leads us to Bayes' formula and to odds.

In Session 2 we meet independent events, and then random variables: numbers
that depend on chance, like the payoff of an investment, with their
expected value and variance.

\say Why does a finance student need this? Banks estimate how likely it is
that a borrower does not repay; models raise alarms; bookmakers quote odds;
an option pays a random amount. All of these are described with exactly the
tools of today.

\begin{cue}
	Today we go slowly. Tell them: ``Before every exercise you will get a few
	minutes to try it yourself, in pairs. Do not wait for me: write something
	down, even if you are not sure.'' Then slide 3 (Session 1 title): just
	click through.
\end{cue}

\subsection{From counting to probability (7 min)}\slide{4}

\say What do we mean when we say that the probability of an event is $0.3$?
The most intuitive answer: if we repeat the experiment many, many times,
the event happens in about 30\% of the repetitions. A probability is a
long-run proportion.

Proportions behave in a very simple way, and the rules of probability
copy that behavior. We write $P(E)$ for the probability of an event $E$.
There are only three basic rules, called the \emph{axioms}.

\board Axioms:
\begin{enumerate}[label=(\arabic*),nosep]
	\item $0\le P(E)\le 1$;
	\item $P(\Om)=1$;
	\item if $E_1,E_2,\dots$ are mutually exclusive (no two of them can
	      happen together), then
	      \[ P\Big(\bigcup_i E_i\Big)=\sum_i P(E_i). \]
\end{enumerate}

\say In words: a probability is between 0 and 1; something always happens,
so the whole sample space has probability 1; and for events that cannot
happen together, the probabilities simply add up.

Everything else follows from these three. Let me show you the most
important consequences, with a tiny example: roll a fair die, so
$\Om=\{1,2,3,4,5,6\}$ and each number has probability $1/6$. Let
$E=\{2,4,6\}$ (``even'') and $F=\{4,5,6\}$ (``at least 4'').

\emph{The complement rule.} $E$ and $E^c$ cannot happen together, and
together they make up $\Om$. So $P(E)+P(E^c)=P(\Om)=1$:
\[ P(E^c)=1-P(E). \]
Example: $P(\text{odd})=1-P(\text{even})=1-\tfrac12=\tfrac12$.

\emph{Monotonicity.} If $E\subseteq F$, then $F$ is $E$ plus something
else, so $P(E)\le P(F)$. A bigger event cannot be less likely.

\emph{The difference rule.} $E$ splits into two mutually exclusive parts:
the part inside $F$, which is $EF$, and the part outside $F$, which is
$EF^c$. So $P(E)=P(EF)+P(EF^c)$, that is,
\[ P(EF^c)=P(E)-P(EF). \]
Example: $EF=\{4,6\}$, so $P(EF^c)=\tfrac12-\tfrac13=\tfrac16$: indeed,
$EF^c=\{2\}$.

\emph{The union rule.} If we add $P(E)$ and $P(F)$, the overlap $EF$ is
counted twice, so we subtract it once:
\[ P(E\cup F)=P(E)+P(F)-P(EF). \]
Example:
$P(E\cup F)=\tfrac12+\tfrac12-\tfrac13=\tfrac23$, and indeed
$E\cup F=\{2,4,5,6\}$ has 4 of the 6 outcomes.

\board Draw the two circles of the slide and label the four regions $EF^c$,
$EF$, $E^cF$, $E^cF^c$.

\say The Venn diagram is the best way to remember all of this. The four
regions are mutually exclusive, and together they make up $\Om$, so their
probabilities add up to 1. Every question about two events is a question
about these four regions.

\ask Which region is ``neither $E$ nor $F$''? \emph{($E^cF^c$, the part
	outside both circles.)} In the die example, which outcomes are there?
\emph{(Only $1$ and $3$: odd and less than 4. Probability $\tfrac26$.)}

\say For three events, the formula is the same inclusion--exclusion that we
used for counts last week: add the singles, subtract the pairs, add back the
triple. It is at the bottom of the slide.

\begin{cue}
	These are the same rules as for counting last week, with probabilities in
	place of counts. If students ask where the rules come from, point out
	that all of them are just the third axiom applied to a split of an event
	into mutually exclusive pieces.
\end{cue}

\subsection{Exercise~\exref{two-exams} --- Two exams (9 min)}\slide{5--6}

\scriptexercise{two-exams}

\begin{cue}
	Only the statement is visible. Give them 3 minutes to try in pairs. Hint
	if they are stuck: ``draw the two circles and start with the overlap''.
\end{cue}

\say Try it in pairs: three minutes. Draw a Venn diagram first.

\begin{cue}
	After 3 minutes: click once (the events and the empty diagram appear).
\end{cue}

\say First we give names to the events. Let $A$ be ``passed Math'' and $B$
``passed Micro''. The data are $P(A)=0.72$, $P(B)=0.64$, and ``both'' is
the intersection: $P(AB)=0.54$.

\say Now we fill the diagram \emph{from the inside out}, exactly as we did
with counts last week.

\begin{center}
	\begin{tikzpicture}[scale=.62]
		\draw[thick] (-2.3,-2) rectangle (4.1,2.2);
		\draw[navy,thick] (0,0) circle (1.5);
		\draw[teal,thick] (1.8,0) circle (1.5);
		\node[navy] at (-1.3,1.55) {$A$};
		\node[teal] at (3.1,1.55) {$B$};
		\node at (-1.95,-1.65) {\small $\Om$};
		\node at (.9,0) {\small $0.54$};
		\node at (-.6,0) {\small $0.18$};
		\node at (2.4,0) {\small $0.10$};
		\node at (3.4,-1.6) {\small $0.18$};
	\end{tikzpicture}
\end{center}

\begin{itemize}
	\item Click: the center, both exams: $0.54$.
	\item Click: Math only. The circle $A$ has probability $0.72$, and $0.54$
	      of it is the overlap, so
	      \[ P(AB^c)=P(A)-P(AB)=0.72-0.54=0.18. \]
	\item Click: Micro only, the same way:
	      \[ P(A^cB)=P(B)-P(AB)=0.64-0.54=0.10. \]
	\item Click: the outside. The four regions add up to 1, so neither exam
	      has probability $1-(0.18+0.54+0.10)=0.18$.
\end{itemize}

\board The diagram with the four numbers. Keep it on the board: we need it
again in Exercise~\exref{exams-conditional}.

\begin{cue}
	Slide 6. The answers appear one by one; for each, ask first, then click.
\end{cue}

\say Now every question is a sum of regions. Let us read the answers off the
diagram and check each one with a rule.

(a) \ask Which regions make up ``at least one exam''? \emph{(The three
	regions inside the circles.)} So $0.18+0.54+0.10=0.82$. With the union rule:
\[ P(A\cup B)=P(A)+P(B)-P(AB)=0.72+0.64-0.54=0.82. \]

(b) Math but not Micro is exactly the region ``Math only'':
$P(AB^c)=0.18$.

(c) \ask How do we write ``neither'' with events? \emph{($A^cB^c$, which
	is the complement of $A\cup B$, by De Morgan.)} So
\[ P(A^cB^c)=1-P(A\cup B)=1-0.82=0.18, \]
the outside region again.

(d) Exactly one: Math only or Micro only. These cannot happen together, so
we add: $0.18+0.10=0.28$. Or: at least one, minus both,
$0.82-0.54=0.28$.

\begin{cue}
	Common mistake: (a) $=0.72+0.64=1.36$. Ask: ``Can a probability be
	larger than 1?'' The students who passed both were counted twice.
\end{cue}

\subsection{Equally likely outcomes (3 min)}\slide{7}

\say In the last exercise the probabilities were given. Where do
probabilities come from? In one important case we can compute them by
counting.

Suppose $\Om$ is finite and every outcome is equally likely. If there are
$|\Om|$ outcomes, each has probability $1/|\Om|$, and an event $E$ with
$|E|$ outcomes has probability
\[ P(E)=\frac{|E|}{|\Om|}=\frac{\text{number of favorable outcomes}}
	{\text{number of all outcomes}}. \]

Tiny example: a fair die, $E$ = ``even'': $3$ favorable outcomes of $6$,
so $P(E)=3/6=1/2$.

\say But careful: ``equally likely'' must come from the setup --- a fair die,
a well-shuffled deck, a person chosen at random. It is not automatic.

\ask We roll two dice and look at the sum. The possible sums are
$2,3,\dots,12$: eleven outcomes. Is the probability of a sum of 7 equal to
$1/11$? \emph{(No. The sums are not equally likely. The 36 ordered pairs
	are: a sum of 7 comes from 6 of them, so its probability is $6/36=1/6$,
	while a sum of 2 comes only from $(1,1)$, probability $1/36$.)}

\say So: choose the sample space so that the outcomes really are equally
likely, then count with the methods of last week. And for ``at least one'',
remember the complement: $P(E)=1-P(E^c)$.

\subsection{Exercise~\exref{five-cards} --- Five cards (8 min)}\slide{8--9}

\scriptexercise{five-cards}

\begin{cue}
	Give them 3 minutes for (a). Hint: ``What is one outcome? How many hands
	are there?'' Some students do not know the deck: 52 cards, 4 suits
	(spades, hearts, diamonds, clubs) of 13 cards each, one ace in each suit.
\end{cue}

\say Before you start: a deck has 52 cards in 4 suits, 13 cards in each
suit, and there are 4 aces, one in each suit. Try part (a), three minutes.

\begin{cue}
	After 3 minutes: click step by step.
\end{cue}

\ask What is one outcome here? Does the order matter? \emph{(A hand of 5
	cards; the order in which the cards arrive does not matter.)}

Click. So the sample space is the set of all 5-card hands:
\[ |\Om|=\binom{52}{5}=2\,598\,960. \]

Click. The deck is well shuffled, so every hand is equally likely, and
$P(E)=|E|/|\Om|$. We only have to count the favorable hands.

(a) Click. Build the hand in two steps. Choose the 2 spades from the 13
spades: $\binom{13}{2}=78$ ways. Click. Choose the 3 diamonds from the 13
diamonds: $\binom{13}{3}=286$ ways. Click. Every choice of spades combines
with every choice of diamonds, so by the basic principle of counting
\[ P(\text{2 spades, 3 diamonds})=\frac{78\cdot286}{2\,598\,960}
	=\frac{22\,308}{2\,598\,960}\approx0.00858. \]
Less than 1\%.

\begin{cue}
	Still on slide 8 (its first line of (b) on slide 9 gives the answer away).
	Give them 2 minutes for (b), with the hint: ``How many cases are there in
	`at least one ace'?'' Then slide 9.
\end{cue}

(b) \ask How many aces can a hand with ``at least one ace'' contain?
\emph{(1, 2, 3 or 4: four cases.)}

\begin{cue}
	Slide 9.
\end{cue}

Four cases to count is a lot of work. Click. The complement is a single case: \emph{no ace}. Click. A hand with
no ace has all 5 cards from the $48$ cards that are not aces:
\[ |\{\text{no ace}\}|=\binom{48}{5}=1\,712\,304. \]
Click. So
\[ P(\text{at least one ace})=1-\frac{1\,712\,304}{2\,598\,960}
	\approx1-0.6588=0.3412. \]

\say Whenever you see ``at least one'', think of the complement first.

\ask A classmate says: choose one ace in 4 ways, then any 4 of the other 51
cards. What is wrong? \emph{(A hand with two aces is counted twice, once for
	each ace as ``the'' ace. The result, about $0.385$, is too large.)}
Click.

\subsection{Conditional probability (6 min)}\slide{10}

\say Now the most important idea of today. Very often we do not know the
outcome, but we know \emph{something} about it. How does that change the
probabilities?

Tiny example.

\ask I roll a fair die and do not show you the result. What is the
probability that it is a 6? \emph{(1/6.)} Now I tell you: ``the number is
even''. What is the probability that it is a 6 now? \emph{(1/3: only
	$2,4,6$ are still possible, and they are equally likely.)}

\say Knowing that $F$ happened, the outcomes outside $F$ are no longer
possible. $F$ becomes the new sample space. Inside $F$ we measure how big
the part of $E$ is that is still possible --- that part is $EF$.

\board Definition: if $P(F)>0$,
\[ P(E\mid F)=\frac{P(EF)}{P(F)}, \]
read ``the probability of $E$ given $F$''.

In the example: $E=\{6\}$, $F=\{2,4,6\}$, $EF=\{6\}$, so
\[ P(E\mid F)=\frac{1/6}{1/2}=\frac13. \]

\say Look at the picture on the slide: $P(E\mid F)$ is the share of the
dark lens $EF$ within the circle $F$.

\say If we multiply the definition by $P(F)$, we get the
\emph{multiplication rule}:
\[ P(EF)=P(F)\,P(E\mid F). \]
It is useful when the conditional probability is what we know. Tiny example:
draw two cards without replacement. The probability that both are aces is
\[ P(\text{1st ace})\,P(\text{2nd ace}\mid\text{1st ace})
	=\frac{4}{52}\cdot\frac{3}{51}=\frac{1}{221}. \]
After one ace is gone, there are 3 aces left among 51 cards. With three
events we keep going: each new factor is conditioned on everything before it,
\[ P(E_1E_2E_3)=P(E_1)\,P(E_2\mid E_1)\,P(E_3\mid E_1E_2). \]

\say Last point on the slide: for a fixed condition $F$, $P(\,\cdot\mid F)$
is a probability itself, so all our rules work for it. For example,
$P(E^c\mid F)=1-P(E\mid F)$.

\begin{cue}
	Stress the order: in $P(E\mid F)$ the condition (what we \emph{know}) is
	on the right, and its probability is in the denominator. Students mix up
	$P(E\mid F)$ and $P(F\mid E)$ all the time; this comes back in Bayes'
	formula.
\end{cue}

\subsection{Exercise~\exref{exams-conditional} --- Two exams, conditionally (8 min)}\slide{11--12}

\scriptexercise{exams-conditional}

\begin{cue}
	The Venn diagram of Exercise~\exref{two-exams} should still be on the
	board. Give them 3 minutes for (a) and (b) in pairs.
\end{cue}

\say These are the same students as in Exercise~\exref{two-exams}, so we can
use our Venn diagram: $0.18$ Math only, $0.54$ both, $0.10$ Micro only,
$0.18$ neither. Try (a) and (b): three minutes. For each part, ask yourself
first: what do I know? That is the condition.

\begin{cue}
	After 3 minutes: click once (the formula and the diagram appear).
\end{cue}

\say The rule: $P(E\mid F)=P(EF)/P(F)$. $F$ happened, so $F$ is the new
sample space.

(a) Click. We know: passed Micro. So the new sample space is the circle
$B$, shaded on the slide, with probability $0.64$. Click. Inside $B$, the
student passed Math in the lens $AB$, $0.54$:
\[ P(A\mid B)=\frac{P(AB)}{P(B)}=\frac{0.54}{0.64}=\frac{27}{32}
	\approx0.844. \]

(b) Click. Now we know: passed Math. The new sample space is the circle $A$,
$0.72$. Click. Inside $A$, ``failed Micro'' is the region Math only,
$0.18$:
\[ P(B^c\mid A)=\frac{P(AB^c)}{P(A)}=\frac{0.18}{0.72}=0.25. \]

\ask Can you get (b) in another way? \emph{(Yes:
	$P(B\mid A)=0.54/0.72=0.75$, and for a fixed condition the complement rule
	works: $1-0.75=0.25$.)}

\begin{cue}
	Common mistake in (a): $0.54/0.72=0.75$. That is $P(B\mid A)$, the answer
	to a different question. The denominator is always the probability of
	what we know.
\end{cue}

\begin{cue}
	Slide 12. Give them 1--2 minutes for (c) first.
\end{cue}

(c) Now we know: passed at least one exam. The new sample space is $A\cup B$,
the two circles together (shaded):
$P(A\cup B)=0.18+0.54+0.10=0.82$.

\ask What is the part of $A\cup B$ where the student passed Math?
\emph{(All of $A$: whoever passed Math passed at least one exam, so
	$A(A\cup B)=A$.)} Click.

Click. So
\[ P(A\mid A\cup B)=\frac{P(A)}{P(A\cup B)}=\frac{0.72}{0.82}
	=\frac{36}{41}\approx0.878. \]

(d) Click. Failed Micro \emph{and} passed at least one exam: then the exam
she passed must be Math, so this is the region Math only, $AB^c$, $0.18$.
Click.
\[ P(B^c\mid A\cup B)=\frac{P(AB^c)}{P(A\cup B)}=\frac{0.18}{0.82}
	=\frac{9}{41}\approx0.220. \]

Click. \say Knowing more changes the probability. Without any information,
the student passed Math with probability $0.72$. If we know she passed
Micro, it is $0.844$; if we know she passed at least one exam, it is
$0.878$.

\subsection{Exercise~\exref{actuarial-exams} --- Three exams in a row (8 min)}\slide{13--14}

\scriptexercise{actuarial-exams}

\say A word on the story: actuaries are the mathematicians of insurance
companies and pension funds, and to become one you pass a series of
professional exams, one after the other.

\begin{cue}
	Give them 3 minutes for (a) and (b). Hint: ``Draw a tree: one branching
	per exam.''
\end{cue}

\say Try (a) and (b), three minutes. Before you compute: are the numbers in
the statement ordinary or conditional probabilities?

\begin{cue}
	After 3 minutes: click step by step.
\end{cue}

Let $E_i$ be ``she passes exam $i$''. \ask Is $0.8$ the probability that she
passes the second exam? \emph{(No: it is the probability that she passes the
	second exam \emph{if} she passed the first, $P(E_2\mid E_1)$.)}
Click. So the data are
\[ P(E_1)=0.9,\qquad P(E_2\mid E_1)=0.8,\qquad P(E_3\mid E_1E_2)=0.7. \]

Click. We draw a tree. From the start, two branches: pass the first exam
($0.9$) or fail it ($0.1$, and she stops). After passing, again two
branches: pass the second ($0.8$) or fail ($0.2$, stop). Then the third:
$0.7$ or $0.3$. The failure probabilities come from the complement rule for
a fixed condition.

\board The tree of the slide, with all six branch probabilities.

Click. \say The probability of a path is the product of the probabilities
along it: that is the multiplication rule.

(a) Click. All three is the top path:
\[ P(E_1E_2E_3)=0.9\cdot0.8\cdot0.7=0.504. \]
In words: 90\% pass the first; 80\% of those pass the second, which is
72\%; 70\% of those pass the third, about half.

\begin{cue}
	Slide 14.
\end{cue}

(b) Fails exactly at the third: passes the first, passes the second, fails
the third. Multiply along that path:
\[ P(E_1E_2E_3^c)=0.9\cdot0.8\cdot0.3=0.216. \]

\begin{cue}
	Give them 2 minutes for (c). Hint: ``What is the condition, and what is
	its probability?''
\end{cue}

(c) Click. The condition is ``not all three'', the complement of the top
path:
\[ P\big((E_1E_2E_3)^c\big)=1-0.504=0.496. \]

Click. ``Failed the second exam'' is the path: passed the first, failed the
second, with probability $0.9\cdot0.2=0.18$. Check with the tree: the three
failure paths are $0.1$, $0.18$ and $0.216$, and they add up to $0.496$.

Click. If she failed the second exam, she certainly did not pass all three,
so the intersection of the two events is $E_1E_2^c$ itself.

Click. So
\[ P\big(E_1E_2^c\mid(E_1E_2E_3)^c\big)=\frac{0.18}{0.496}
	=\frac{45}{124}\approx0.363. \]

\ask Why not just $P(E_2^c)$? \emph{(If she fails the first exam, she stops,
	so she does not pass the second one either; $P(E_2^c)=0.1+0.18=0.28$ would
	include her. ``Failed the second exam'' means she took it and failed.)}

\subsection{Total probability and Bayes' formula (5 min)}\slide{15}

\say The tree of the last exercise is the key to the next two formulas.

Tiny example. On 30\% of the days it rains. On a rainy day my bus is late
with probability $0.5$; on a dry day only with probability $0.1$. Let $F$
be ``it rains'' and $E$ be ``the bus is late''. So $P(F)=0.3$,
$P(E\mid F)=0.5$, $P(E\mid F^c)=0.1$.

\board The tree of the slide: first $F$ or $F^c$, then $E$ or $E^c$, with
the numbers of the example.

\emph{First question: what is the probability that the bus is late?} The bus
can be late on a rainy day or on a dry day: two paths of the tree that end
in $E$, and they cannot happen together. Multiply along each path, then add:
\[ P(E)=P(E\mid F)P(F)+P(E\mid F^c)P(F^c)=0.5\cdot0.3+0.1\cdot0.7
	=0.15+0.07=0.22. \]
This is the \emph{law of total probability}: split according to $F$ or
$F^c$, multiply along the paths, add the paths that end in $E$.

\emph{Second question: the bus was late. What is the probability that it
	rained?} Now we know $E$ and ask about $F$: we need $P(F\mid E)$, but the
data give $P(E\mid F)$. By the definition, $P(F\mid E)=P(EF)/P(E)$. The
numerator is the path ``rain, then late'', $0.15$; the denominator is all
paths that end in ``late'', $0.22$:
\[ P(F\mid E)=\frac{0.15}{0.22}=\frac{15}{22}\approx0.68. \]

\board Bayes' formula:
\[ P(F\mid E)=\frac{P(E\mid F)P(F)}{P(E\mid F)P(F)+P(E\mid F^c)P(F^c)}. \]

\say Bayes' formula turns the condition around. In the tree: one path,
divided by all the paths that end in $E$. Note what happened: rain has
probability $0.3$, but once we see the late bus, it goes up to about $0.68$.
The evidence changed our belief.

\subsection{Exercise~\exref{early-warning} --- An early-warning model, finance (11 min)}\slide{16--17}

\scriptexercise{early-warning}

\say Some words on the story. A borrower \emph{defaults} when she does not
repay her loan as agreed. Banks use \emph{early-warning models}: statistical
models that raise a flag on borrowers who look likely to default, so that
the bank can act in time. No model is perfect: it misses some defaulters,
and it also flags some good borrowers.

\ask Before any computation: a borrower is flagged. Make a quick guess ---
what is the probability that she defaults? \emph{(Collect a few guesses.
	Many will say about 90\%. Write the guesses on the board; we come back to
	them.)}

\begin{cue}
	Give them 4 minutes in pairs for (a) and (b). Hint: ``Name the events,
	write the three numbers as probabilities, draw the tree.'' After 4
	minutes: click step by step.
\end{cue}

Let $D$ be ``defaults'' and $A$ ``flagged'' ($A$ for alarm).
\ask Which number is $P(A\mid D)$, and which is $P(D\mid A)$?
\emph{(``Flags 90\% of the defaulters'': among defaulters, 90\% are flagged,
	so $P(A\mid D)=0.9$. And $P(A\mid D^c)=0.05$. $P(D\mid A)$ is what (b) asks
	for.)}

Click. So $P(D)=0.02$, $P(A\mid D)=0.9$, $P(A\mid D^c)=0.05$. On the tree: first
default or not ($0.02$ and $0.98$), then flag or no flag. Multiply along the
paths:
\[ P(DA)=0.02\cdot0.9=0.018,\qquad P(DA^c)=0.02\cdot0.1=0.002, \]
\[ P(D^cA)=0.98\cdot0.05=0.049,\qquad P(D^cA^c)=0.98\cdot0.95=0.931. \]

\begin{center}
	\begin{tikzpicture}[x=1.2cm,y=.62cm,lab/.style={font=\small,inner sep=1pt}]
		\node (r) at (0,0) {$\bullet$};
		\node (d) at (1.3,1.5) {$D$};
		\node (g) at (1.3,-1.5) {$D^c$};
		\draw (r) -- node[lab,above left] {$0.02$} (d);
		\draw (r) -- node[lab,below left] {$0.98$} (g);
		\node (da) at (2.6,2.25) {$A$};
		\node (dn) at (2.6,.75) {$A^c$};
		\node (ga) at (2.6,-.75) {$A$};
		\node (gn) at (2.6,-2.25) {$A^c$};
		\draw (d) -- node[lab,above left] {$0.9$} (da);
		\draw (d) -- node[lab,below left] {$0.1$} (dn);
		\draw (g) -- node[lab,above left] {$0.05$} (ga);
		\draw (g) -- node[lab,below left] {$0.95$} (gn);
		\node[lab,right,text=teal] at (2.8,2.25) {$\mathbf{0.018}$};
		\node[lab,right] at (2.8,.75) {$0.002$};
		\node[lab,right,text=teal] at (2.8,-.75) {$\mathbf{0.049}$};
		\node[lab,right] at (2.8,-2.25) {$0.931$};
	\end{tikzpicture}
\end{center}

(a) Click. A flagged borrower is either a defaulter or a good borrower: two
paths ending in $A$. Total probability (click once more for the numbers):
\[ P(A)=P(A\mid D)P(D)+P(A\mid D^c)P(D^c)=0.018+0.049=0.067. \]
So $6.7\%$ of all borrowers are flagged.

(b) Click. Bayes: the path ``defaults and flagged'' divided by all paths
ending in a flag:
\[ P(D\mid A)=\frac{P(A\mid D)P(D)}{P(A)}=\frac{0.018}{0.067}=\frac{18}{67}
	\approx0.269. \]

\say Only about 27\%! Compare that with your guesses.

\begin{cue}
	Slide 17. The counts make the surprise understandable; go through the
	tree of 10\,000 borrowers slowly, number by number.
\end{cue}

\say Let us check this with counts, which is much easier to believe. Imagine
$10\,000$ borrowers. 2\% of them, $200$, will default; $9800$ will not. The
model flags 90\% of the 200 defaulters: $180$, and misses $20$. It flags 5\%
of the 9800 good borrowers: $490$, and does not flag $9310$.

Click. So $180+490=670$ borrowers are flagged, and only $180$ of them default:
\[ P(D\mid A)=\frac{180}{670}\approx0.269. \]

(c) Click. Not flagged: $20+9310=9330$ borrowers, of whom $20$ default:
\[ P(D\mid A^c)=\frac{20}{9330}=\frac{2}{933}\approx0.00214. \]

Click. \say This is the \emph{base-rate} effect. The model is good: it
catches 90\% of the defaulters. Yet only about 27\% of the flags are true.
Why? Defaults are rare. The 5\% false flags come from a huge group of good
borrowers, so they outnumber the true flags: $490$ against $180$.

Click. \say But the model is still very useful: a flag raises the
probability of default from $2\%$ to $27\%$, more than thirteen times, and
no flag lowers it to about $0.2\%$.

\begin{cue}
	Common mistake: answering (b) with $0.9$. That is $P(A\mid D)$; the
	question is $P(D\mid A)$. Turning the condition around is exactly what
	Bayes' formula does.
\end{cue}

\begin{bg}{default, early-warning systems and the base rate}
	\emph{Default}: the borrower does not pay interest or principal as agreed
	(in bank regulation, typically 90 days past due, or ``unlikely to pay'').
	The \emph{probability of default (PD)} over one year is a key number of
	credit risk; under the Basel rules banks estimate it for every borrower,
	and it drives their capital requirements and loan pricing. A 2\% annual
	PD is realistic for an average retail or small-business portfolio.

	\emph{Early-warning system (EWS)}: a model, often a scorecard or a
	machine-learning classifier on account data (late payments, overdraft use,
	falling salary inflow), that flags clients before they default. A flag
	triggers action: a call from the relationship manager, a restructuring of
	the loan, no new credit lines.

	In statistics language: $P(A\mid D)=0.9$ is the \emph{sensitivity}
	(hit rate, true positive rate), $P(A\mid D^c)=0.05$ is the
	\emph{false positive rate} (false alarm rate), and $P(D\mid A)\approx0.27$
	is the \emph{precision} (positive predictive value). The 2\% is the
	\emph{base rate} (prevalence). With a low base rate, precision is low even
	for a good model --- the same effect as in medical screening tests.

	\terms default; probability of default (PD); early-warning system;
	sensitivity / hit rate; false positive rate; precision; base rate.

	\sq{Why not tune the model so that it flags fewer good borrowers?}
	\ans One can raise the threshold, but then it also catches fewer
	defaulters: there is always a trade-off between missed defaults and false
	alarms (in statistics, the ROC curve). Which mix is best depends on the
	costs: a missed default costs the loan, a false alarm only a phone call.

	\sq{If 73\% of the flags are false, is the model useless?}
	\ans No. The bank can afford to look at $6.7\%$ of its clients
	closely, and that group contains $90\%$ of the future defaults. Without
	the model it would have to check everyone.
\end{bg}

\subsection{Odds (4 min)}\slide{18}

\say There is another way of describing how likely something is, which you
know from betting: \emph{odds}. The odds of an event $H$ are the probability
that it happens divided by the probability that it does not:
\[ \text{odds of }H=\frac{P(H)}{P(H^c)}=\frac{P(H)}{1-P(H)}. \]

Tiny examples. If $P(H)=2/3$, the odds are $\frac{2/3}{1/3}=2$: ``2 to 1 in
favor'' --- two cases where $H$ happens for every one where it does not. If
$P(H)=0.2$, the odds are $0.2/0.8=1/4$: ``4 to 1 against''.

\say And back: odds $o$ mean $o$ ``yes'' cases for every 1 ``no'' case, so
out of $o+1$ cases, $o$ are ``yes'':
\[ P(H)=\frac{o}{1+o}. \]
Check: odds $2$ give $2/3$.

\board
\[ \text{odds}=\frac{P(H)}{1-P(H)},\qquad P(H)=\frac{\text{odds}}{1+\text{odds}}. \]

\say Why are odds useful? Because Bayes' formula becomes very simple. Write
Bayes' formula for $H$ and for $H^c$, given evidence $E$, and divide the two:
the long denominators are the same and cancel. What remains is
\[ \frac{P(H\mid E)}{P(H^c\mid E)}=\frac{P(H)}{P(H^c)}\cdot
	\frac{P(E\mid H)}{P(E\mid H^c)}. \]
New odds equal old odds times the \emph{likelihood ratio}: how many times
more likely the evidence is if $H$ is true than if it is false.

\say Back to the bus. Old odds of rain: $0.3/0.7=3/7$. A late bus is
$0.5/0.1=5$ times as likely on a rainy day. New odds: $\frac37\cdot5=\frac{15}{7}$,
so $P(\text{rain}\mid\text{late})=\frac{15}{15+7}=\frac{15}{22}$ --- the same
as before, with less work.

\say So: evidence that is more likely if $H$ is true raises the odds of $H$;
evidence that is less likely if $H$ is true lowers them.

\subsection{Exercise~\exref{early-warning-odds} --- Two warning signals: odds, finance (8 min)}\slide{19--20}

\scriptexercise{early-warning-odds}

\begin{cue}
	Give them 3 minutes for (a). Hint: ``prior odds, then likelihood ratio,
	then multiply, then convert back to a probability.''
\end{cue}

\say This is the same bank and the same model. Try part (a) with odds: three
minutes.

\begin{cue}
	After 3 minutes: click step by step.
\end{cue}

Click. The rule: new odds = old odds $\times$ likelihood ratio, and odds $o$
correspond to the probability $o/(1+o)$.

Click. Prior odds of default, before any signal:
\[ \frac{P(D)}{P(D^c)}=\frac{0.02}{0.98}=\frac{1}{49}. \]
One defaulter for every 49 good borrowers.

Click. Likelihood ratio of a flag:
\[ \frac{P(A_1\mid D)}{P(A_1\mid D^c)}=\frac{0.90}{0.05}=18. \]
A flag is 18 times as likely for a defaulter as for a good borrower.

Click. (a) Posterior odds: $\frac{1}{49}\cdot18=\frac{18}{49}$, so
\[ P(D\mid A_1)=\frac{18}{18+49}=\frac{18}{67}\approx0.269. \]

Click. \say The same as in Exercise~\exref{early-warning}: among the flagged,
there are 18 defaulters for every 49 good borrowers --- $180$ against $490$
in our 10\,000 borrowers.

\begin{cue}
	Slide 20. Explain the second model slowly; conditional independence is
	only needed in this informal form (independence comes properly in
	Session 2).
\end{cue}

\say Now a second model, with the same rates. The statement says that its
errors are independent of the first model's errors, given whether the
borrower defaults. In plain words: among defaulters, whether model 2 flags
someone has nothing to do with whether model 1 did; and the same among good
borrowers. Then the likelihood ratios simply \emph{multiply}.

(b) Click. Both models flag: the likelihood ratio is $18\cdot18=324$, so
\[ \text{odds}=\frac{1}{49}\cdot18\cdot18=\frac{324}{49},\qquad
	P(D\mid A_1A_2)=\frac{324}{324+49}=\frac{324}{373}\approx0.869. \]

(c) \ask The second model does \emph{not} flag her. What is the likelihood
ratio of ``no flag''? \emph{(Use the complements: $P(A_2^c\mid D)=0.10$ and
	$P(A_2^c\mid D^c)=0.95$.)} Click.
\[ \frac{P(A_2^c\mid D)}{P(A_2^c\mid D^c)}=\frac{0.10}{0.95}=\frac{2}{19}. \]
It is less than 1: a missing flag lowers the odds.

Click.
\[ \text{odds}=\frac{1}{49}\cdot18\cdot\frac{2}{19}=\frac{36}{931},\qquad
	P(D\mid A_1A_2^c)=\frac{36}{36+931}=\frac{36}{967}\approx0.0372. \]

Click. \say Two flags make default likely, about 87\%. One flag that the
other model does not confirm leaves only about 4\%, not far above the 2\%
we started from.

\begin{cue}
	Common mistake: forgetting to convert the odds back. Odds of $324/49$
	are larger than 1; they are not a probability.
\end{cue}

\begin{bg}{odds, likelihood ratios and credit scorecards}
	\emph{Betting odds.} Bookmakers quote ``fractional'' odds (UK, e.g.\ 4/1
	means a win of 4 per 1 staked, i.e.\ implied probability $1/5$) or
	``decimal'' odds (continental Europe, e.g.\ 5.0 means the total payout
	per unit staked, again implied probability $1/5$). The implied
	probabilities of all outcomes add up to slightly more than 1; the excess
	is the bookmaker's margin.

	\emph{Log-odds and scorecards.} Taking logarithms turns the odds form of
	Bayes' formula into a sum: log-odds after the evidence $=$ log-odds before
	$+$ log of the likelihood ratio. This is why credit \emph{scorecards}
	add up points: each characteristic of the client contributes points, and
	the total score is (a linear function of) the log-odds of being a good
	payer. A typical convention: every 20 points double the odds. Logistic
	regression, the workhorse of credit scoring, models exactly the log-odds.

	\emph{Conditional independence.} Multiplying the likelihood ratios is
	correct only if the signals are independent given $D$ and given $D^c$ (the
	``naive Bayes'' assumption). Two models built on the same account data
	tend to make the same mistakes; then a second flag adds much less
	information than the factor 18 suggests, and the 87\% is too optimistic.

	\terms odds (fractional, decimal); implied probability; likelihood ratio;
	log-odds; scorecard; logistic regression; conditional independence;
	naive Bayes.

	\sq{Why do banks use odds instead of probabilities?}
	\ans Because evidence acts multiplicatively on odds (additively on
	log-odds), so independent pieces of information can be combined one by
	one, in any order. With probabilities you would have to redo the whole
	Bayes formula every time.

	\sq{Is it realistic that the two models' errors are independent?}
	\ans Only roughly, and only if the models use different information (for
	example, one uses account data and the other market data or the loan
	officer's judgment). That is why banks like signals from different
	sources.
\end{bg}

\subsection{Summary (3 min)}\slide{21}

\begin{itemize}
	\item $P(E^c)=1-P(E)$; $P(E\cup F)=P(E)+P(F)-P(EF)$;
	      $P(EF^c)=P(E)-P(EF)$. Fill the Venn diagram from the inside out.
	\item Equally likely outcomes: $P(E)=|E|/|\Om|$; ``at least one''
	      $\to$ complement.
	\item $P(E\mid F)=P(EF)/P(F)$: $F$ is the new sample space. Multiplication
	      rule: multiply along a path of a tree.
	\item Total probability: add the paths that end in $E$. Bayes' formula:
	      one path divided by all of them.
	\item Odds: new odds $=$ old odds $\times$ likelihood ratio.
	\item Rare events and imperfect signals: most alarms are false.
\end{itemize}

\say If you remember one picture from this session, let it be the tree:
multiply along the paths, add the paths you need, and divide when you know
something.

\say Every exercise of today is in the study guide on the course website,
with a detailed, step-by-step solution. The Exercises page there has more
problems for self-study.

\begin{cue}
	End of Session 1. Announce the exact restart time after the short break
	(e.g.\ ``we continue at \ldots, in ten minutes''), and leave the tree of
	the early-warning model on the board: early in Session 2 comes the
	generalized Bayes formula, which is the same tree with more branches.
\end{cue}

%======================================================================
\section{Session 2: Independence and random variables}
%======================================================================

\subsection{Recap of Session 1 (4 min)}\slide{22--23}

\begin{cue}
	Slide 22 (Session 2 title) --- just click through. Slide 23 is the recap.
	Let the students settle; ask the recap questions and really wait for the
	answers.
\end{cue}

Welcome back. Before we go on, let's collect the tools of the first
session, because we will use every one of them in the next ninety minutes.

\say Probabilities of combined events: the union and the complement.
\[ P(E\cup F)=P(E)+P(F)-P(EF),\qquad P(E^c)=1-P(E). \]

\ask Why do we subtract $P(EF)$ in the first formula? \emph{(The outcomes in
	both events are counted once in $P(E)$ and once more in $P(F)$.)}

\say Conditional probability: if we know that $F$ has happened, $F$ becomes
the new sample space, and only the part $EF$ of $E$ counts:
\[ P(E\mid F)=\frac{P(EF)}{P(F)},\qquad\text{so}\qquad P(EF)=P(F)\,P(E\mid F). \]
The second form is the multiplication rule: in a tree we multiply along a
path.

\say With two hypotheses, $F$ and $F^c$, we drew a tree. The law of total
probability adds the paths that end in $E$; Bayes' formula takes one of these
paths and divides it by all of them. And in odds form: new odds equal old
odds times the likelihood ratio.

\ask In the early-warning model of Exercise~\exref{early-warning}, what was
the surprising lesson? \emph{(When the event is rare and the signal is
	imperfect, most alarms are false alarms.)}

\say Now three new things. First, more than two hypotheses --- the same tree,
just with more branches. Second, independent events: when does knowing one
event not change the chances of another? And third, random variables:
numbers that depend on chance, with their distribution, expected value and
standard deviation. That last part is where finance really starts: payoffs,
profits and losses are random variables.

\subsection{The generalized Bayes formula (5 min)}\slide{24}

\say In Session 1 we had two hypotheses: $F$ or $F^c$. Often there are more:
a loan can be in one of three rating classes, an economy can be in boom,
normal times or recession. What we need is a \emph{partition}.

\say $F_1,\dots,F_n$ is a \emph{partition} of $\Om$ if the events are
mutually exclusive and exactly one of them occurs.

Tiny example: roll a die, and let $F_1=\{1,2\}$, $F_2=\{3,4\}$,
$F_3=\{5,6\}$. Whatever we roll, exactly one of these three events occurs.
$F$ and $F^c$ form the simplest partition, with $n=2$.

\ask Is $\{1,2,3\}$, $\{3,4,5\}$, $\{6\}$ a partition? \emph{(No: the first
	two overlap, a 3 is in both.)}

\say If $E$ is any event, then $E$ is cut into pieces by the partition:
$EF_1$, $EF_2$, \dots, $EF_n$. These pieces are mutually exclusive, so their
probabilities add up, and each piece is a path of the tree:
$P(EF_i)=P(E\mid F_i)\,P(F_i)$.

\board
\[ P(E)=\sum_{i=1}^n P(E\mid F_i)\,P(F_i),\qquad
	P(F_j\mid E)=\frac{P(E\mid F_j)\,P(F_j)}{\sum_{i=1}^n P(E\mid F_i)\,P(F_i)}. \]

\say The first is the law of total probability, the second the generalized
Bayes formula. In words, exactly as before: add the paths that end in $E$;
and Bayes is one path divided by all of them. The only change is that the
tree now has $n$ branches instead of 2.

Tiny example, on the board. We have three boxes. Box 1 contains only red
balls, box 2 is half red and half blue, box 3 contains only blue balls. We
pick a box at random and draw a ball; let $R$ be ``the ball is red''.
\[ P(R)=\tfrac13\cdot1+\tfrac13\cdot\tfrac12+\tfrac13\cdot0=\tfrac12. \]
The ball is red. Which box did it come from?
\[ P(\text{box }1\mid R)=\frac{1/3}{1/2}=\frac23,\qquad
	P(\text{box }2\mid R)=\frac{1/6}{1/2}=\frac13,\qquad
	P(\text{box }3\mid R)=0. \]

\ask Why is box 1 now twice as likely as box 2? \emph{(Box 1 produces red
	balls twice as often; before the draw, all three boxes were equally
	likely.)}

\begin{cue}
	Point at the tree on the slide: one path to $E$ for each $F_i$. The
	posterior probabilities $P(F_j\mid E)$ always add up to 1; tell them to
	use this as a check.
\end{cue}

\subsection{Exercise~\exref{default-by-rating} --- Defaults by rating class, finance (10 min)}\slide{25--27}

\scriptexercise{default-by-rating}

\begin{cue}
	Show only the statement and give them three minutes to name the events,
	draw the tree and try (a). Walk around; most students write $P(D)=0.08$
	somewhere --- that is the trap of part (b).
\end{cue}

\ask What are the hypotheses here, and what is the evidence? \emph{(The
	hypotheses are the rating classes; the evidence is default or no
	default.)}

\say Name the events first. $A$, $B$, $C$: the loan is rated A, B, C. $D$:
the loan defaults.

\say $A$, $B$, $C$ form a partition: every loan is in exactly one class.

\board The data, sorted:
\[ \begin{aligned}
		 & P(A)=0.5,       &  & P(B)=0.3,       &  & P(C)=0.2;        \\
		 & P(D\mid A)=0.01, &  & P(D\mid B)=0.03, &  & P(D\mid C)=0.08.
	\end{aligned} \]

\say Notice that the default probabilities are \emph{conditional}: $0.08$ is
the chance of default for a loan rated C, not for any loan.

Now the tree: the first level is the rating class, with $0.5$, $0.3$,
$0.2$. The second level is default or not. The ``no default'' branches are
the complements: $0.99$, $0.97$, $0.92$.

\board Draw the tree; at the end of every path write the product of the
probabilities along it: $0.005$, $0.495$; $0.009$, $0.291$; $0.016$,
$0.184$.

\ask What should the six numbers at the end add up to? \emph{(To 1: they
	are all the possible outcomes.)}

(a) \say Next slide. A path: multiply. For example
$P(DA)=P(D\mid A)P(A)=0.5\cdot0.01=0.005$. The three paths that end in $D$
are mutually exclusive, so we add them:
\[ P(D)=0.5\cdot0.01+0.3\cdot0.03+0.2\cdot0.08=0.005+0.009+0.016=0.03. \]
So 3\% of the loans default. This is a weighted average of the three default
rates, with the class shares as weights --- that is why it lies between
$0.01$ and $0.08$.

(b) \say Bayes: one path ending in $D$, divided by all of $P(D)$.
\[ P(A\mid D)=\frac{P(D\mid A)P(A)}{P(D)}=\frac{0.005}{0.03}=\frac16\approx0.167. \]
\ask Who can do the other two? \emph{(B: $0.009/0.03=0.3$; C:
	$0.016/0.03=8/15\approx0.533$.)}
\[ P(B\mid D)=\frac{0.009}{0.03}=0.3,\qquad
	P(C\mid D)=\frac{0.016}{0.03}=\frac{8}{15}\approx0.533. \]
Check: $0.167+0.3+0.533=1$. A defaulted loan was rated A, B or C.

\begin{cue}
	Common mistake: answering $0.08$ for $P(C\mid D)$. Write the two numbers
	next to each other on the board: $P(D\mid C)=0.08$ is the default rate of
	a C loan, $P(C\mid D)\approx0.533$ is the share of C loans among the
	defaulted ones.
\end{cue}

(c) \say Last slide of the exercise. Same idea, but now the paths that end in
$D^c$. First the denominator, by the complement:
$P(D^c)=1-0.03=0.97$. The path through $C$ to $D^c$:
$P(D^c\mid C)P(C)=0.92\cdot0.2=0.184$. So
\[ P(C\mid D^c)=\frac{P(D^c\mid C)P(C)}{P(D^c)}=\frac{0.184}{0.97}\approx0.190. \]

\ask Why is this so close to the prior $P(C)=0.2$? \emph{(Almost every loan
	survives the year, so ``no default'' tells us very little.)}

\say The lesson: class C is only 20\% of the loans, but it gives more than
half, 53\%, of the defaults. Class A is half of the loans, but gives only one
sixth of the defaults.

\begin{bg}{internal ratings and probabilities of default}
	\emph{Internal rating}: a bank's own grade for every borrower, based on
	financial statements, payment history, industry, etc. For every grade the
	bank estimates a \emph{probability of default} (PD), the chance that the
	borrower stops paying within one year. These are exactly the numbers
	$P(D\mid A)$, $P(D\mid B)$, $P(D\mid C)$ of the exercise. The portfolio
	default rate $P(D)=0.03$ is the PD of a randomly chosen loan: a weighted
	average of the class PDs.

	\emph{Why the Bayes direction matters.} The PDs answer ``how risky is
	this class?''. Bayes answers the opposite question: ``where do our
	defaults come from?'' --- important for the bank's credit committee (is
	the loss concentrated in one class?) and for checking the rating system:
	a good rating system puts most of the later defaults in the worst
	classes. Here class C gives $8/15\approx53\%$ of the defaults, so the
	ratings clearly separate good and bad borrowers.

	\emph{What the PDs are used for.} The \emph{expected loss} of a loan is
	$\text{EL}=\text{PD}\cdot\text{LGD}\cdot\text{EAD}$: probability of
	default times \emph{loss given default} (the share that is not recovered)
	times \emph{exposure at default} (the amount owed). The bank prices this
	into the interest rate; under the Basel rules the PDs also determine how
	much capital it must hold against unexpected losses.

	\terms internal rating; probability of default (PD); loss given default
	(LGD); exposure at default (EAD); expected loss; Basel capital rules.

	\sq{Is a 3\% default rate realistic?}
	\ans For a portfolio of loans to smaller, riskier companies, yes; for
	large, highly rated companies the one-year default rates are far below
	1\%. Default rates also rise sharply in recessions, so a bank's PDs are
	estimated over a whole economic cycle.

	\sq{If C loans are so bad, why does the bank lend to them at all?}
	\ans Because it charges them a higher interest rate. If the extra
	interest covers the expected loss and the cost of capital, the C loans
	can even be the most profitable part of the portfolio.
\end{bg}

\subsection{Independent events (6 min)}\slide{28}

\say Usually, knowing that $F$ has happened changes the probability of $E$:
that was the whole point of conditional probability. Sometimes, however, it
does not change anything. Then we call the events independent.

\say $E$ and $F$ are \emph{independent} if $P(EF)=P(E)\,P(F)$.

Why this formula? If $P(F)>0$, divide both sides by $P(F)$:
\[ P(E\mid F)=\frac{P(EF)}{P(F)}=P(E). \]
So knowing $F$ does not change the chances of $E$. We use the product form
as the definition because it is symmetric in $E$ and $F$ and it also works
when $P(F)=0$.

Tiny example. Roll a fair die. Let $E=\{2,4,6\}$ (even) and $F=\{1,2\}$ (at
most 2). Then $EF=\{2\}$, and
\[ P(EF)=\tfrac16,\qquad P(E)P(F)=\tfrac12\cdot\tfrac13=\tfrac16. \]
Equal: $E$ and $F$ are independent. Indeed, among $\{1,2\}$ half of the
numbers are even, just as among all six numbers.

\ask And with $G=\{1,2,3\}$ (at most 3)? \emph{($EG=\{2\}$, $P(EG)=1/6$,
	but $P(E)P(G)=1/2\cdot1/2=1/4$: not independent. Among 1, 2, 3 only one
	number is even.)}

\say Independence is a property of the numbers, not of the story: you check
the equation.

If $E$ and $F$ are independent, then $E$ and $F^c$ are independent too ---
if knowing that $F$ happened tells us nothing about $E$, then knowing that
$F$ did not happen tells us nothing either.

\say Three events $E$, $F$, $G$ are independent if \emph{all four} product
rules hold:
\[ P(EF)=P(E)P(F),\ \ P(EG)=P(E)P(G),\ \ P(FG)=P(F)P(G),\ \
	P(EFG)=P(E)P(F)P(G). \]
Doesn't the last one follow from the first three? No --- the next exercise
gives an example where it fails.

\ask Two events are mutually exclusive, each with positive probability. Can
they be independent? \emph{(No: if one occurs, the other cannot, so
	$P(EF)=0$, while $P(E)P(F)>0$.)}

\begin{cue}
	Students very often confuse ``mutually exclusive'' with ``independent''.
	Say it explicitly: mutually exclusive events are strongly dependent ---
	knowing one tells you the other did not happen.
\end{cue}

\say The most useful consequence: $n$ independent trials, each a success
with probability $p$. ``At least one success'' is the complement of ``no
success at all'', and for independent trials we multiply:
\[ P(\text{at least one success})=1-(1-p)^n. \]
For example, three tosses of a fair coin: $P(\text{at least one head})=
1-(1/2)^3=7/8$.

\subsection{Exercise~\exref{pairwise-not-mutual} --- Pairwise but not mutually independent (7 min)}\slide{29--30}

\scriptexercise{pairwise-not-mutual}

\begin{cue}
	Give them two minutes for (a): list $\Om$ and the three events as sets,
	then check the three product rules. Most students guess that $A$ and $C$
	are dependent; let them find out that they are not.
\end{cue}

\say Four equally likely outcomes, each with probability $\frac14$:
\[ \Om=\{HH,HT,TH,TT\}. \]
The first letter is the first coin.

\ask Write $A$, $B$, $C$ as sets. \emph{($A=\{HH,HT\}$, $B=\{HH,TH\}$,
	$C=\{HH,TT\}$.)}

Each contains two of the four outcomes, so $P(A)=P(B)=P(C)=\frac12$.

\board The table from the slide: rows are the outcomes, a dot means the
event of the column occurs.
\begin{center}
	\begin{tabular}{c|ccc}
		     & $A$       & $B$       & $C$       \\ \hline
		$HH$ & $\bullet$ & $\bullet$ & $\bullet$ \\
		$HT$ & $\bullet$ &           &           \\
		$TH$ &           & $\bullet$ &           \\
		$TT$ &           &           & $\bullet$
	\end{tabular}
\end{center}

(a) \say Look at the table: the only outcome that is in two columns at once
is $HH$. So every pair has the same intersection, $AB=AC=BC=\{HH\}$, with
probability $\frac14$. And every product of two probabilities is
$\frac12\cdot\frac12=\frac14$. So
\[ P(AB)=P(A)P(B),\qquad P(AC)=P(A)P(C),\qquad P(BC)=P(B)P(C). \]
Every pair is independent: the events are \emph{pairwise independent}.

\ask $A$ and $B$ are clearly independent: two separate coins. But why are
$A$ and $C$ independent? \emph{(If the first coin is heads, ``same side''
	means the second coin is heads too, which still has probability $\frac12$:
	$P(C\mid A)=\frac12=P(C)$.)}

(b) \say Next slide. For three events we need all four equations. The first
three hold by (a). The fourth: all three events occur only at $HH$, so
$ABC=\{HH\}$, and
\[ P(ABC)=\tfrac14,\qquad\text{but}\qquad
	P(A)P(B)P(C)=\tfrac12\cdot\tfrac12\cdot\tfrac12=\tfrac18. \]
$\frac14\ne\frac18$, so $A$, $B$, $C$ are \emph{not} independent.

(c) \say If $A$ and $B$ both occur, the outcome is $HH$, so $C$ occurs for
sure:
\[ P(C\mid AB)=1\ne\tfrac12=P(C). \]
One coin alone says nothing about ``same side''; two of the events together
decide the third.

\say The moral: pairwise independence is not enough. For three events, check
all four equations.

\subsection{Exercise~\exref{bond-defaults} --- Bond defaults: independent or not?, finance (9 min)}\slide{31--33}

\scriptexercise{bond-defaults}

\begin{cue}
	Give them three minutes for (a) and (b). Before they start, make sure
	they translate the two swaps into events (the first reveal of the slide).
\end{cue}

\ask Which event makes the first-to-default swap pay, and which the
second-to-default swap? \emph{(At least one default: $D_1\cup D_2$. Both
	default: $D_1D_2$.)}

\say For the union we always have inclusion--exclusion,
$P(D_1\cup D_2)=P(D_1)+P(D_2)-P(D_1D_2)$. So in every case the key number
is $P(D_1D_2)$.

(a) \say Independent: the probability of both is the product.
\[ P(D_1D_2)=P(D_1)P(D_2)=0.05\cdot0.05=0.0025, \]
\[ P(D_1\cup D_2)=0.05+0.05-0.0025=0.0975. \]

\say Check with the complement. ``No default'' is $D_1^cD_2^c$ (De Morgan),
and the complements of independent events are independent, each with
probability $0.95$:
\[ P(D_1\cup D_2)=1-0.95^2=1-0.9025=0.0975. \]

(b) \say Next slide. Now the two issuers depend on the same economy, and
$P(D_1D_2)=0.02$ is given.

\ask Can we multiply $0.05\cdot0.05$ here? \emph{(No: the product rule holds
	only for independent events.)}

\say $0.02$ is $8$ times as large as $0.0025$. The union is again
inclusion--exclusion:
\[ P(D_1\cup D_2)=0.05+0.05-0.02=0.08. \]

(c) \say The definition of conditional probability:
\[ P(D_2\mid D_1)=\frac{P(D_1D_2)}{P(D_1)}=\frac{0.02}{0.05}=0.4. \]
$0.4\ne0.05=P(D_2)$: the defaults are \emph{not independent}. Once bond 1
has defaulted --- probably because the economy is in trouble --- bond 2
defaults with probability $0.4$ instead of $0.05$. This is \emph{positive}
dependence.

(d) \say Last slide. Back to independence, but now with 10 bonds.

\ask How would you compute ``at least one default'' for 10 bonds?
\emph{(Not with inclusion--exclusion --- far too many terms. Complement: at
	least one default is the opposite of ``all 10 bonds survive''.)}

Each bond survives with probability $0.95$; the bonds are independent, so we
multiply:
\[ P(\text{no default})=0.95^{10}\approx0.599,\qquad
	P(\text{at least one default})=1-0.95^{10}\approx0.401. \]
This is the formula $1-(1-p)^n$ from the theory slide.

\say Each bond is fairly safe, but with 10 of them a first-to-default swap
pays with probability about 40\%.

\board The comparison table:
\begin{center}
	\begin{tabular}{lcc}
		two bonds    & at least one & both     \\ \hline
		independent  & $0.0975$     & $0.0025$ \\
		same economy & $0.08$       & $0.02$
	\end{tabular}
\end{center}

\say Lesson: positive dependence makes ``at least one'' slightly less likely,
but ``both'' much more likely. The defaults come together, so they overlap
more.

\begin{bg}{default correlation and basket swaps}
	\emph{First-to-default (FtD) and second-to-default (StD) swaps}: insurance
	contracts on a basket of bonds. The protection buyer pays a regular
	premium (the \emph{spread}); the FtD pays at the first default in the
	basket, the StD only at the second. In this exercise the FtD pays on
	$D_1\cup D_2$, the StD on $D_1D_2$.

	\emph{Default correlation.} Companies depend on the same economy: in a
	recession many of them get into trouble at the same time. This common
	(\emph{systematic}) factor makes defaults positively dependent. One
	measures it by the correlation of the default indicators,
	\[ \rho=\frac{P(D_1D_2)-P(D_1)P(D_2)}{\sqrt{P(D_1)(1-P(D_1))\,
				P(D_2)(1-P(D_2))}}
		=\frac{0.02-0.0025}{0.05\cdot0.95}\approx0.37 \]
	in case (b); under independence $\rho=0$. (Students will meet correlation
	later; you do not need it in class.)

	\emph{Why it matters for prices.} The table of the last slide is the
	whole story: with more correlation the StD event becomes much more likely
	(its seller wants a higher spread), while the FtD event becomes slightly
	less likely. Banks model this with factor models (e.g.\ the Gaussian
	copula). In 2008 the correlation of mortgage defaults was badly
	underestimated: products that only lose after many defaults (the senior
	tranches of mortgage CDOs) lost much more than the models predicted.

	\terms first-to-default / second-to-default swap; spread; default
	correlation; systematic risk; factor model; Gaussian copula.

	\sq{Where would a number like $P(D_1D_2)=0.02$ come from?}
	\ans From history (how often did similar companies default in the same
	year?) or from market prices of basket products, from which the
	correlation the market assumes can be backed out. Either way it is an
	estimate, and a very uncertain one.

	\sq{Is a diversified bond portfolio then safe?}
	\ans Diversification removes the company-specific part of the risk, but
	not the common factor. In a deep recession many bonds default together,
	whatever the number of bonds.
\end{bg}

\subsection{Discrete random variables (6 min)}\slide{34}

\say Until now we asked whether an event happens. In finance we usually ask
\emph{how much}: how much do we win, how much does the portfolio lose. For
that we need random variables.

\say A \emph{random variable} $X$ assigns a number to every outcome. It is
\emph{discrete} if it takes finitely or countably many values.

Tiny example. Toss two fair coins, $\Om=\{HH,HT,TH,TT\}$, and let $X$ be the
number of heads. Then $X(HH)=2$, $X(HT)=X(TH)=1$, $X(TT)=0$. So $X$ is not
random in a mysterious way: it is just a function on $\Om$, and the
randomness comes from the outcome.

\say The \emph{probability mass function}, pmf for short, gives the
probability of each value: $p(x)=P(X=x)$. It is never negative, and the
values add up to 1.

\board For the two coins:
\begin{center}
	\begin{tabular}{c|ccc}
		$x$    & 0             & 1             & 2             \\\hline
		$p(x)$ & $\frac14$     & $\frac12$     & $\frac14$
	\end{tabular}
\end{center}

\ask Why is $p(1)=\frac12$? \emph{(Two outcomes, $HT$ and $TH$, give one
	head: $\frac24$.)}

\say The \emph{distribution function} adds up the pmf from the left:
\[ F(x)=P(X\le x)=\sum_{y\le x}p(y). \]
For the two coins: $F(x)=0$ for $x<0$, then $\frac14$ from $0$, then
$\frac14+\frac12=\frac34$ from $1$, and $1$ from $2$ on.

\say So $F$ is a step function. It jumps by $p(x)$ at every value $x$, it is
constant in between, and it rises from 0 to 1.

\begin{cue}
	Point at the graph on the slide (there $p(1)=\frac14$, $p(2)=\frac12$,
	$p(3)=p(4)=\frac18$). Stress the dots: at a jump the filled dot is on the
	\emph{upper} step, because $F(x)=P(X\le x)$ already includes $p(x)$. This
	is the most common mistake in the sketch of the urn exercise.
\end{cue}

\ask On the slide's graph, what is $F(2)$, and what is $F(2.5)$?
\emph{(Both $\frac34$: $F(2)$ includes $p(2)$, and there is no value between
	2 and 2.5.)}

\say From $F$ we can read off interval probabilities:
\[ P(a<X\le b)=F(b)-F(a),\qquad P(X>a)=1-F(a). \]
For example, on the slide $P(1<X\le3)=F(3)-F(1)=\frac78-\frac14=\frac58$,
which is $p(2)+p(3)=\frac12+\frac18$. Watch the inequalities: $a$ itself is
not included, $b$ is.

\subsection{Exercise~\exref{urn-game} --- A game with an urn (12 min)}\slide{35--37}

\scriptexercise{urn-game}

\begin{cue}
	The longest exercise of the day: let them work in pairs on (a) for four
	minutes. Hint if they are stuck: ``list the possible color pairs first''.
	Then go through the slides reveal by reveal.
\end{cue}

\ask How many balls are in the urn, and how many pairs can we draw?
\emph{($8+4+2=14$ balls; the order does not matter, so $\binom{14}{2}=91$
	pairs.)}

(a) \say All 91 pairs are equally likely. We list the possible pairs by their
colors and count each kind.
\begin{itemize}
	\item two white: $\binom82=28$;
	\item white and orange: $8\cdot2=16$ (one of the 8 white balls and one of
	      the 2 orange ones);
	\item two orange: $1$;
	\item white and black: $8\cdot4=32$;
	\item black and orange: $4\cdot2=8$;
	\item two black: $\binom42=6$.
\end{itemize}
Check: $28+16+1+32+8+6=91$, so no pair is missing.

\ask What is $X$ for each color pair? \emph{(WW: $-2$; WO: $-1$; OO: $0$; WB:
	$-1+2=1$; BO: $2$; BB: $4$.)}

\say Different color pairs give different values, so each probability is the
number of pairs divided by 91:
\begin{center}
	\renewcommand{\arraystretch}{1.3}
	\begin{tabular}{c|cccccc}
		$x$    & $-2$            & $-1$            & $0$            & $1$
		       & $2$             & $4$                                                \\\hline
		$p(x)$ & $\frac{28}{91}$ & $\frac{16}{91}$ & $\frac{1}{91}$ & $\frac{32}{91}$
		       & $\frac{8}{91}$  & $\frac{6}{91}$
	\end{tabular}
\end{center}
Check: the probabilities add up to $\frac{91}{91}=1$.

(b) \say Next slide. Add up the pmf from the left: $28$, $28+16=44$,
$44+1=45$, $45+32=77$, $77+8=85$, $85+6=91$, all over 91.

\board
\[ F(x)=\begin{cases}
		0,             & x<-2,       \\
		28/91,         & -2\le x<-1, \\
		44/91,         & -1\le x<0,  \\
		45/91,         & 0\le x<1,   \\
		77/91,         & 1\le x<2,   \\
		85/91,         & 2\le x<4,   \\
		1,             & x\ge 4.
	\end{cases} \]
and sketch the steps: filled dot at the left end of every step, open circle
at the right end.

\ask What is $F(-2)$? \emph{($\frac{28}{91}$: ``$X\le-2$'' includes
	$X=-2$.)}

(c) \say $X>0$ is the complement of $X\le0$:
\[ P(X>0)=1-F(0)=1-\frac{45}{91}=\frac{46}{91}\approx0.505. \]
Directly: the positive values are 1, 2 and 4, so
$\frac{32+8+6}{91}=\frac{46}{91}$. And
\[ F(1)=P(X\le1)=\frac{77}{91}=\frac{11}{13}\approx0.846. \]

\begin{cue}
	If someone answers $F(1)=\frac{45}{91}$: that is $P(X<1)$. The value 1
	itself belongs to $F(1)$.
\end{cue}

(d) \say Last slide. We have not defined the expected value yet --- it comes
properly on the next theory slide --- but the idea is natural: each value
times its probability, then add.
\[ \E[X]=\frac{(-2)\cdot28+(-1)\cdot16+0\cdot1+1\cdot32+2\cdot8+4\cdot6}{91}
	=\frac{-56-16+0+32+16+24}{91}=0. \]
On average we neither win nor lose: a \emph{fair game}.

\say The variance is the average squared distance of the values from the
mean. The shortcut is $\Var(X)=\E[X^2]-(\E[X])^2$. For $\E[X^2]$, square the
values and keep the probabilities:
\[ \E[X^2]=\frac{4\cdot28+1\cdot16+0\cdot1+1\cdot32+4\cdot8+16\cdot6}{91}
	=\frac{112+16+0+32+32+96}{91}=\frac{288}{91}. \]
\say The variance is in squared euros; its square root, the standard
deviation $\SD(X)$, is back in euros.
\[ \Var(X)=\frac{288}{91}-0^2=\frac{288}{91}\approx3.165,\qquad
	\SD(X)=\sqrt{\frac{288}{91}}\approx1.779\text{ euros}. \]

\say The game is fair, but risky: a typical deviation from 0 is about 1.78
euros.

\begin{cue}
	If you run late, do (d) only for $\E[X]$ and leave the variance for the
	next slide, where the formula is introduced.
\end{cue}

\subsection{Expected value, variance, standard deviation (6 min)}\slide{38}

\say The \emph{expected value} of $X$ is the average of its values, weighted
by their probabilities:
\[ \E[X]=\sum_x x\,p(x). \]
Think of it as the long-run average: if we play the game many times, the
average winnings per game get close to $\E[X]$.

Tiny example, the two coins again: $X$ is the number of heads, with
probabilities $\frac14,\frac12,\frac14$ for $0,1,2$:
\[ \E[X]=0\cdot\tfrac14+1\cdot\tfrac12+2\cdot\tfrac14=1. \]

\say A function of $X$: apply the function to every value, keep the
probabilities.
\[ \E[g(X)]=\sum_x g(x)\,p(x). \]
For example $\E[X^2]=0\cdot\frac14+1\cdot\frac12+4\cdot\frac14=\frac32$.
Note: $\frac32$ is not $(\E[X])^2=1$. In general $\E[g(X)]$ is not
$g(\E[X])$; we see a finance example of this in a minute.

\say There is one important exception: linear functions.
$\E[aX+b]=a\,\E[X]+b$. If every head pays 10 euros and the game costs 5
euros, the net gain is $10X-5$, and its expected value is
$10\cdot1-5=5$ euros.

\say The expected value tells us the center; it says nothing about the
spread. The \emph{variance} is the expected squared distance from the mean
$\mu=\E[X]$:
\[ \Var(X)=\E\big[(X-\mu)^2\big]=\E[X^2]-\big(\E[X]\big)^2. \]

For the coins, from the definition: $(0-1)^2\cdot\frac14+(1-1)^2\cdot\frac12
	+(2-1)^2\cdot\frac14=\frac12$. With the shortcut: $\frac32-1^2=\frac12$. The
same number; the shortcut is usually faster.

\say The variance is in squared units --- squared euros --- so we take the
square root: the \emph{standard deviation}, $\SD(X)=\sqrt{\Var(X)}$. For the
coins $\SD(X)=\sqrt{1/2}\approx0.707$.

\board
\[ \Var(aX+b)=a^2\Var(X),\qquad \SD(aX+b)=|a|\,\SD(X). \]

\ask Why does $b$ disappear? \emph{(Adding a constant shifts every value and
	the mean by the same amount; the distances from the mean do not change.)}
\ask Why $a^2$? \emph{(Every distance from the mean is multiplied by $a$, and
	the variance squares the distances.)}

So for $10X-5$: $\SD=10\cdot0.707\approx7.07$ euros.

\say In finance: $\E[X]$ is the expected return or payoff, and $\SD(X)$
measures the \emph{risk}.

\subsection{Exercise~\exref{call-option} --- The payoff of a call option, finance (9 min)}\slide{39--40}

\scriptexercise{call-option}

\begin{cue}
	Explain the call option in two sentences before they start (the first
	paragraph below). Then give them three minutes for (a).
\end{cue}

\say A quick word on the finance. A \emph{call option} with strike price 110
gives you the right, but not the obligation, to buy the stock for 110 in one
year. If the stock then costs more than 110, you buy at 110 and gain the
difference; if it costs less, you simply do not use the option. So the payoff
is $C=\max(S-110,0)$: never negative.

(a) \say Each price times its probability, then add.
\[ \E[S]=80\cdot0.2+100\cdot0.4+120\cdot0.3+140\cdot0.1=16+40+36+14=106. \]
For the variance, first the mean of $S^2$:
\[ \E[S^2]=6400\cdot0.2+10000\cdot0.4+14400\cdot0.3+19600\cdot0.1
	=1280+4000+4320+1960=11560. \]
\[ \Var(S)=\E[S^2]-(\E[S])^2=11560-106^2=11560-11236=324,\qquad
	\SD(S)=\sqrt{324}=18. \]
The price is typically about 18 away from 106.

(b) \say Next slide. $C$ is a function of $S$: compute the payoff for each
price.

\ask What does the option pay if $S=80$? $100$? $120$? $140$? \emph{(0, 0,
	10, 30.)}

\say Two prices give $C=0$, so their probabilities add up:
\[ p_C(0)=0.2+0.4=0.6,\qquad p_C(10)=0.3,\qquad p_C(30)=0.1. \]
\[ \E[C]=0\cdot0.6+10\cdot0.3+30\cdot0.1=3+3=6. \]
\[ \E[C^2]=0+100\cdot0.3+900\cdot0.1=120,\qquad \Var(C)=120-6^2=84,\qquad
	\SD(C)=\sqrt{84}\approx9.17. \]

(c) \ask What do we get if we put the expected price into the payoff?
\emph{($\max(106-110,0)=0$.)}

\say But $\E[C]=6$. The expected price is below the strike, so ``plugging in
the average'' says the option is worthless. But the price is random: with
probability $0.4$ it ends above 110, and then the option pays; otherwise it
pays 0, never a negative amount. So the average payoff is positive.

\say In general $\E[g(X)]\ne g(\E[X])$. Apply $g$ first, then average:
$\E[g(X)]=\sum_x g(x)\,p(x)$.

\begin{bg}{call options and why uncertainty has value}
	\emph{Call option}: the right, but not the obligation, to buy an asset
	(the \emph{underlying}) at the \emph{strike price} $K$ on the
	\emph{expiry date}. Its payoff $\max(S-K,0)$ has the famous
	``hockey stick'' graph: zero below $K$, then rising one-for-one. A
	\emph{put} is the right to sell, with payoff $\max(K-S,0)$. The buyer
	pays a price up front (the \emph{premium}). An option with $S$ above $K$
	is \emph{in the money}, below $K$ \emph{out of the money}.

	\emph{The lesson of part (c)} is a central idea of option theory: since
	the payoff function is convex (it bends upward), $\E[\max(S-K,0)]\ge
	\max(\E[S]-K,0)$ (Jensen's inequality). The option keeps the upside and
	cuts off the downside, so more uncertainty in $S$ means a more valuable
	option. That is why option prices rise with \emph{volatility}, the
	standard deviation of returns.

	\terms underlying; strike; expiry; premium; in/out of the money;
	payoff diagram; volatility; convexity.

	\sq{So is 6 the price of this option?}
	\ans Not exactly. The fair price today discounts the payoff (money in a
	year is worth less than money today), and in option pricing theory the
	expectation is taken with special ``risk-neutral'' probabilities, not
	with the real-world probabilities of the exercise. But the structure ---
	price = (discounted) expected payoff --- is exactly this computation.

	\sq{Why is $\SD(C)\approx9.17$ smaller than $\SD(S)=18$, if options are
		``risky''?}
	\ans In euros the payoff varies less, but the option costs much less than
	the stock. Relative to the money invested, an option is far riskier: the
	payoff is 0 with probability $0.6$, so the buyer loses the whole premium
	more often than not.
\end{bg}

\subsection{Exercise~\exref{trading-strategy} --- A simple trading strategy, finance (7 min)}\slide{41--42}

\scriptexercise{trading-strategy}

\begin{cue}
	Two minutes for (a); then ask (b) and (c) before revealing them --- they
	are one-line answers with the rules of slide 38.
\end{cue}

(a) \say The values of $X$: $300$ with probability $0.55$, and $-200$ with
probability $0.45$. A loss is a negative number.
\[ \E[X]=300\cdot0.55+(-200)\cdot0.45=165-90=75. \]
On average the strategy earns 75 euros a day.

For the variance, $\E[X^2]$ first:
\[ \E[X^2]=300^2\cdot0.55+(-200)^2\cdot0.45=49\,500+18\,000=67\,500. \]
\[ \Var(X)=\E[X^2]-\big(\E[X]\big)^2=67\,500-75^2=67\,500-5625=61\,875. \]

\say Next slide. $61\,875$ is in squared euros; back to euros:
\[ \SD(X)=\sqrt{61\,875}\approx248.7\text{ euros}. \]
\say The typical deviation, about 249 euros, is much larger than the mean, 75
euros: on a single day the result is very uncertain.

\board Rules for $aX+b$: $\E[aX+b]=a\E[X]+b$, \ $\SD(aX+b)=|a|\,\SD(X)$.

(b) \ask The position is doubled: what are $a$ and $b$? \emph{($a=2$,
	$b=0$.)}
\[ \E[2X]=2\cdot75=150,\qquad \SD(2X)=2\cdot\SD(X)\approx497.5. \]
\say Leverage doubles the expected profit, but it doubles the risk too.

\begin{cue}
	Common mistake: $\Var(2X)=2\Var(X)$. The variance is multiplied by $4$
	(it is $247\,500$); only the standard deviation is doubled.
\end{cue}

(c) \ask A fixed fee of 20 euros: $a$ and $b$? \emph{($a=1$, $b=-20$.)}
\[ \E[X-20]=75-20=55,\qquad \SD(X-20)=\SD(X)\approx248.7. \]
\say A constant shifts every outcome by the same amount --- $280$ and $-220$
instead of $300$ and $-200$. The center moves, the spread does not change.

\begin{bg}{P\&L, leverage and risk-adjusted return}
	\emph{P\&L} (profit and loss): the money a trading account gains or loses
	over a period. \emph{Position size}: how much is invested; doubling it
	doubles every gain and every loss. \emph{Leverage}: taking a larger
	position with borrowed money. \emph{Volatility}: the standard deviation
	of the P\&L (or of returns), the standard measure of risk.

	\emph{Risk-adjusted return.} Traders compare strategies by the ratio of
	mean to standard deviation (the idea behind the \emph{Sharpe ratio}).
	Here it is $75/248.7\approx0.30$ per day. Leverage does not change it:
	$150/497.5\approx0.30$ --- both numbers double. The fee lowers it:
	$55/248.7\approx0.22$. So leverage scales a strategy up or down, but it
	does not make it better; costs make it worse.

	\terms P\&L; position size; leverage; volatility; Sharpe ratio;
	transaction costs; margin call.

	\sq{If the expected profit doubles, why not use as much leverage as
		possible?}
	\ans Because losses double too. A run of bad days can wipe out the
	account, and the broker demands more collateral (a \emph{margin call})
	exactly when the trader is losing. Many trading firms failed this way.

	\sq{Is a 20-euro fee a big deal when the mean is 75?}
	\ans It takes more than a quarter of the expected daily profit. For
	active strategies, transaction costs often decide whether the strategy
	is profitable at all.
\end{bg}

\subsection{Summary (3 min)}\slide{43}

\begin{itemize}
	\item Partition $F_1,\dots,F_n$: $P(E)=\sum_i P(E\mid F_i)P(F_i)$, and
	      the generalized Bayes formula gives $P(F_j\mid E)$: one path divided
	      by all of them.
	\item Independence: $P(EF)=P(E)P(F)$; for three events all four product
	      rules; pairwise independence is not enough. ``At least one'' in $n$
	      independent trials: $1-(1-p)^n$.
	\item pmf $p(x)=P(X=x)$; distribution function $F(x)=P(X\le x)$, a step
	      function.
	\item $\E[X]=\sum x\,p(x)$; $\E[g(X)]=\sum g(x)p(x)$, in general not
	      $g(\E[X])$.
	\item $\Var(X)=\E[X^2]-(\E[X])^2$; $\Var(aX+b)=a^2\Var(X)$;
	      $\SD(X)=\sqrt{\Var(X)}$.
\end{itemize}

\say That was a long day: from probabilities of events, through conditional
probability and Bayes, to independence and random variables. Next week we
continue with random variables. Every exercise of today is in the study
guide on the course website, with a detailed, step-by-step solution, and the
Exercises page there has more problems on every topic of today for
self-study. Thank you, and see you next week.

\begin{cue}
	If there is reserve time, let them start one of the self-study exercises
	on random variables in class and walk around. Next week: start with a
	two-minute recap of the pmf, $\E[X]$ and $\Var(X)$, using the urn game of
	Exercise~\exref{urn-game}.
\end{cue}


\end{document}
